\documentclass[11pt,a4paper]{article}

% ── Codificación y fuentes ────────────────────────────────────────────────────
\usepackage[utf8]{inputenc}
\usepackage[T1]{fontenc}
\usepackage{lmodern}
\usepackage[spanish]{babel}

% ── Geometría y espaciado ─────────────────────────────────────────────────────
\usepackage[top=2.5cm, bottom=2.5cm, left=2.8cm, right=2.8cm]{geometry}
\usepackage{setspace}
\setstretch{1.15}
\usepackage{parskip}

% ── Tablas ────────────────────────────────────────────────────────────────────
\usepackage{booktabs}
\usepackage{tabularx}
\usepackage{array}
\usepackage{longtable}
\usepackage{enumitem}

% ── Colores y cajas ───────────────────────────────────────────────────────────
\usepackage[dvipsnames]{xcolor}
\usepackage{tcolorbox}
\tcbuselibrary{skins, breakable}

% ── Cabeceras y pies de página ────────────────────────────────────────────────
\usepackage{fancyhdr}
\usepackage{lastpage}
\pagestyle{fancy}
\fancyhf{}
\fancyhead[L]{\small\textbf{Informe de Evaluación} --- Física para Bioanalistas}
\fancyhead[R]{\small thamairys sosa 33685285}
\fancyfoot[C]{\small Página \thepage\ de \pageref{LastPage}}
\renewcommand{\headrulewidth}{0.4pt}

% ── Hiperenlaces ──────────────────────────────────────────────────────────────
\usepackage[colorlinks=true, linkcolor=NavyBlue, urlcolor=NavyBlue]{hyperref}

% ── Paleta de colores personalizados ─────────────────────────────────────────
\definecolor{desempeno_excelente}{HTML}{1A6B2F}
\definecolor{desempeno_bueno}{HTML}{2E5A9C}
\definecolor{desempeno_suficiente}{HTML}{8B6914}
\definecolor{desempeno_deficiente}{HTML}{C0392B}
\definecolor{desempeno_insuficiente}{HTML}{7B241C}
\definecolor{grisFondo}{HTML}{F5F5F5}
\definecolor{azulTitulo}{HTML}{1B2A6B}
\definecolor{verdeFortaleza}{HTML}{1A6B2F}
\definecolor{naranjaMejora}{HTML}{A04000}

% ── Estilos de cajas ─────────────────────────────────────────────────────────
\tcbset{
  cajaTitulo/.style={
    enhanced, breakable,
    colback=azulTitulo!8, colframe=azulTitulo,
    fonttitle=\bfseries\large, coltitle=white,
    attach boxed title to top left={yshift=-2mm, xshift=4mm},
    boxed title style={colback=azulTitulo},
    top=6pt, bottom=6pt, left=8pt, right=8pt,
  },
  cajaFortaleza/.style={
    enhanced, breakable,
    colback=verdeFortaleza!6, colframe=verdeFortaleza!60,
    fonttitle=\bfseries, coltitle=verdeFortaleza,
    top=4pt, bottom=4pt, left=8pt, right=8pt,
  },
  cajaMejora/.style={
    enhanced, breakable,
    colback=naranjaMejora!6, colframe=naranjaMejora!60,
    fonttitle=\bfseries, coltitle=naranjaMejora,
    top=4pt, bottom=4pt, left=8pt, right=8pt,
  },
  cajaRecomendacion/.style={
    enhanced, breakable,
    colback=grisFondo, colframe=gray!50,
    fonttitle=\bfseries, coltitle=black,
    top=4pt, bottom=4pt, left=8pt, right=8pt,
  },
}

% ─────────────────────────────────────────────────────────────────────────────
\begin{document}

% ══ PORTADA ══════════════════════════════════════════════════════════════════
\begin{center}
  {\Large\bfseries\color{azulTitulo} Universidad de Oriente/Núcleo Sucre --- Física para Bioanalistas}\\[4pt]
  {\large Informe de Evaluación Preliminar}\\[12pt]
  \rule{\linewidth}{1.2pt}\\[6pt]
  {\LARGE\bfseries thamairys sosa 33685285}\\[4pt]
  \rule{\linewidth}{0.6pt}
\end{center}

% ══ FICHA TÉCNICA ════════════════════════════════════════════════════════════
\vspace{4pt}
\begin{tcolorbox}[cajaTitulo, title=Datos del informe]
\begin{tabular}{@{}llll@{}}
  \textbf{Estudiante:}  & thamairys sosa 33685285  &
  \textbf{Fecha:}       & 2026-06-26 \\[4pt]
  \textbf{Archivo PDF:} & \texttt{thamairys\_sosa\_33685285.pdf}  &
  \textbf{Páginas:}     & 4 \\
\end{tabular}
\end{tcolorbox}

% ══ RESULTADO GLOBAL ═════════════════════════════════════════════════════════
\vspace{6pt}
\begin{center}
  \begin{tcolorbox}[
    enhanced,
    colback=white, colframe=desempeno_bueno,
    width=0.62\linewidth,
    boxrule=2pt,
    halign=center, valign=center,
    top=8pt, bottom=8pt,
  ]
    \centering
    {\large\bfseries Puntaje sugerido}\\[6pt]
    {\Huge\bfseries\color{desempeno_bueno} 7.9/10}\\[6pt]
    {\large Nivel de desempeño: \textbf{\color{desempeno_bueno} Bueno}}
  \end{tcolorbox}
\end{center}

% ══ RESUMEN GENERAL ══════════════════════════════════════════════════════════
\section*{Resumen general}
El estudiante demuestra una sólida comprensión de los principios fundamentales de la mecánica de fluidos, aplicando correctamente las fórmulas y metodologías para la mayoría de los problemas. Se observa un buen manejo de las ecuaciones de continuidad, Bernoulli, Poiseuille y el principio de Arquímedes. Sin embargo, se detectaron errores recurrentes de índole procedimental, principalmente en cálculos aritméticos, manejo de exponentes y una interpretación ligeramente incompleta de la pregunta en un caso, lo que afectó la precisión de los resultados numéricos.

% ══ FORTALEZAS Y ÁREAS DE MEJORA ════════════════════════════════════════════
\vspace{4pt}
\begin{minipage}[t]{0.48\linewidth}
  \begin{tcolorbox}[cajaFortaleza, title={Fortalezas}]
\begin{itemize}[leftmargin=1.5em, itemsep=2pt]
  \item Correcta identificación y aplicación de los principios físicos y fórmulas relevantes (Presión hidrostática, Continuidad, Bernoulli, Arquímedes, Poiseuille).
  \item Buen entendimiento de las conversiones de unidades (mm a m).
  \item Presentación clara y organizada de los pasos de resolución en la mayoría de los problemas.
  \item Correcta manipulación algebraica para el despeje de variables.
\end{itemize}
  \end{tcolorbox}
\end{minipage}
\hfill
\begin{minipage}[t]{0.48\linewidth}
  \begin{tcolorbox}[cajaMejora, title={Areas de mejora}]
\begin{itemize}[leftmargin=1.5em, itemsep=2pt]
  \item Precisión en los cálculos aritméticos y el manejo de potencias y exponentes.
  \item Revisión cuidadosa de los cálculos iniciales a partir de los datos del problema (ej. peso a partir de masa).
  \item Lectura e interpretación exhaustiva de la pregunta para asegurar que la respuesta final aborde el requerimiento exacto (ej. 'cambio porcentual' vs. 'porcentaje remanente').
\end{itemize}
  \end{tcolorbox}
\end{minipage}

% ══ OBSERVACIONES POR PREGUNTA ═══════════════════════════════════════════════
\section*{Observaciones por pregunta}

\begin{longtable}{>{\bfseries}p{2.8cm} p{1.6cm} p{7.0cm} p{2.8cm}}
  \toprule
  \textbf{Pregunta} & \textbf{Puntaje} & \textbf{Evaluación} & \textbf{Errores detectados} \\
  \midrule
  \endhead
  \bottomrule
  \endfoot
    \textbf{Pregunta 1: Presión Hidrostática y Venosa} & 2.0/2.0 & La respuesta es completamente correcta. El estudiante identificó correctamente la presión relevante, aplicó la fórmula de presión hidrostática y realizó los cálculos de manera precisa, incluyendo la conversión de unidades y la expresión del resultado final. & \textit{---} \\
    \midrule
    \textbf{Pregunta 2: Principio de Arquímedes} & 1.2/2.0 & El estudiante aplicó correctamente el principio de Arquímedes para calcular la fuerza de empuje, el volumen del objeto y su densidad. Sin embargo, cometió un error inicial al calcular el peso del objeto en el aire (W\_aire) a partir de la masa dada (0.50 kg * 9.8 m/s$^{2}$ = 4.9 N, no 4.50 N). Este error se propagó a los cálculos subsiguientes de la fuerza de empuje, el volumen y la densidad del tejido. & \textit{Error en el cálculo inicial del peso en el aire (W\_aire).; Propagación de error en el cálculo de la fuerza de empuje.; Propagación de error en el cálculo del volumen del tejido.; Propagación de error en el cálculo de la densidad del tejido.} \\
    \midrule
    \textbf{Pregunta 3: Hemodinámica (Continuidad y Bernoulli)} & 2.0/2.0 & La respuesta es completamente correcta. El estudiante aplicó de manera impecable las ecuaciones de continuidad y Bernoulli, realizó los despejes y sustituciones correctamente, y obtuvo el resultado numérico exacto para la presión en la zona estrecha. & \textit{---} \\
    \midrule
    \textbf{Pregunta 4: Ley de Poiseuille (Análisis Proporcional)} & 1.5/2.0 & El estudiante comprendió correctamente la relación de la Ley de Poiseuille con el radio (Q $\propto$ r\textasciicircum{}4) y calculó correctamente la proporción del nuevo flujo respecto al original (Q2 = 0.4096 Q1). Sin embargo, la pregunta solicitaba el "porcentaje de diferencia" o "cambio porcentual", y la respuesta final solo indica el porcentaje del flujo original que permanece (40.96\%). Faltó calcular el cambio porcentual real, que sería una reducción del 59.04\%. & \textit{Interpretación incompleta de la pregunta sobre 'porcentaje de diferencia' o 'cambio porcentual'.} \\
    \midrule
    \textbf{Pregunta 5: Flujo Viscoso y Resistencia} & 1.2/2.0 & El estudiante identificó y aplicó correctamente la Ley de Poiseuille para calcular la diferencia de presión. Los datos fueron correctamente identificados y convertidos. No obstante, se cometió un error procedimental significativo en el cálculo de r\textasciicircum{}4 en el denominador ((4 x 10$^{-}$$^{4}$ m)$^{4}$ no es 2.56 x 10$^{-}$$^{1}$$^{3}$, sino 256 x 10$^{-}$$^{1}$$^{6}$ o 2.56 x 10$^{-}$$^{1}$$^{4}$). Este error en el manejo de potencias y exponentes afectó el resultado final. & \textit{Error en el cálculo de la potencia del radio (r\textasciicircum{}4).; Error en el manejo de exponentes.} \\
    \midrule
\end{longtable}

% ══ ERRORES DETECTADOS ═══════════════════════════════════════════════════════
\section*{Análisis de errores}

\subsection*{Errores conceptuales}
\begin{itemize}[leftmargin=1.5em, itemsep=2pt]
  \item No se detectaron errores conceptuales fundamentales en la comprensión de los principios físicos.
\end{itemize}

\subsection*{Errores procedimentales}
\begin{itemize}[leftmargin=1.5em, itemsep=2pt]
  \item Cálculo incorrecto del peso en el aire a partir de la masa (Pregunta 2).
  \item Omisión del cálculo final del cambio porcentual solicitado (Pregunta 4).
  \item Error en el cálculo de una potencia con exponentes (r\textasciicircum{}4) (Pregunta 5).
  \item Errores de cálculo aritmético que se propagan a resultados finales.
\end{itemize}

% ══ RECOMENDACIONES AL ESTUDIANTE ════════════════════════════════════════════
\section*{Retroalimentación para el estudiante}
\begin{tcolorbox}[cajaRecomendacion, title={Dirigido a: thamairys sosa 33685285}]
Thomairys, tu examen demuestra una sólida comprensión de los principios fundamentales de la mecánica de fluidos. Has aplicado correctamente la mayoría de las fórmulas y has seguido los pasos lógicos para resolver los problemas. Sin embargo, es crucial que prestes más atención a los detalles en los cálculos numéricos, especialmente al manejar potencias y exponentes, y al realizar cálculos iniciales a partir de los datos del problema. Revisa cuidadosamente cada paso aritmético y asegúrate de que tu respuesta final aborde exactamente lo que se pregunta. Practicar más con la calculadora y verificar las unidades en cada paso te ayudará a mejorar la precisión.
\end{tcolorbox}

% ══ NOTA PARA EL PROFESOR ════════════════════════════════════════════════════
\section*{Nota para el profesor}
\begin{tcolorbox}[
  enhanced, breakable,
  colback=yellow!8, colframe=yellow!60!black,
  fonttitle=\bfseries, coltitle=black,
  title={Observaciones para revisión manual},
  top=4pt, bottom=4pt, left=8pt, right=8pt,
]
El estudiante demuestra un buen dominio conceptual de la unidad de fluidos. Los errores principales son de naturaleza procedimental (cálculos aritméticos, manejo de exponentes) y una ligera falta de precisión en la interpretación de la pregunta 4 (calculó el porcentaje remanente en lugar del cambio porcentual). No se observan errores conceptuales graves. El puntaje sugerido refleja el entendimiento de los principios a pesar de los errores de cálculo.
\end{tcolorbox}

% ══ PIE DE INFORME ════════════════════════════════════════════════════════════
\vspace{12pt}
\begin{center}
  \footnotesize\color{gray}
  Informe generado automáticamente por el sistema de evaluación con IA (gemini-2.5-flash) y soporte LaTex. \\
  Este documento es una evaluación \textbf{preliminar} para ser revisado por el profesor antes de ser entregado al 
  estudiante. \\
  Generado el 2026-06-26 \textbullet\ BIOANALISIS Sistema Académico de Evaluación.
  \end{center}

\end{document}
